\documentclass[a4paper,14pt,oneside,openany]{memoir}

%%% Поля и отступы %%%
\usepackage[left=30mm, right=15mm, top=20mm, bottom=20mm]{geometry}
\pagestyle{plain}
\parindent=1.25cm
\usepackage{indentfirst}

%%% Языковые параметры и шрифт %%%
\usepackage{polyglossia}
\setdefaultlanguage[indentfirst=true,forceheadingpunctuation=false]{russian}
\setotherlanguages{english}

\IfFontExistsTF{Times New Roman}
{\setmainfont{Times New Roman}
\newfontfamily\cyrillicfont{Times New Roman}[Script=Cyrillic]}
{\setmainfont[Script=Cyrillic]{FreeSerif}}

\IfFontExistsTF{Courier New}
{\setmonofont{Courier New}
\newfontfamily\cyrillicfonttt{Courier New}[Script=Cyrillic]}
{\setmonofont[Script=Cyrillic]{FreeMono}}

\IfFontExistsTF{Arial}
{\setsansfont{Arial}
\newfontfamily\cyrillicfontsf{Arial}[Script=Cyrillic]}
{\setsansfont[Script=Cyrillic]{FreeSans}}

%%% Стиль заголовков %%%
\setsecnumdepth{subsection}
\renewcommand*{\chapterheadstart}{}
\renewcommand*{\printchaptername}{}
\renewcommand*{\chapnumfont}{\normalfont\bfseries}
\renewcommand*{\afterchapternum}{\hspace{1em}}
\renewcommand*{\printchaptertitle}{\normalfont\bfseries\centering\MakeUppercase}
\setbeforesecskip{20pt}
\setaftersecskip{20pt}
\setsecheadstyle{\raggedright\normalfont\bfseries}
\setbeforesubsecskip{20pt}
\setaftersubsecskip{20pt}
\setsubsecheadstyle{\raggedright\normalfont\bfseries}

%%% Оглавление %%%
\addto\captionsrussian{\renewcommand\contentsname{Содержание}}
\setrmarg{2.55em plus1fil}
\renewcommand{\aftertoctitle}{\afterchaptertitle \vspace{-\cftbeforechapterskip}}
\renewcommand*{\cftchapternumwidth}{1.5em}
\renewcommand*{\cftchapterfont}{\normalfont\MakeUppercase}
\renewcommand*{\cftchapterpagefont}{\normalfont}
\renewcommand*{\cftchapterdotsep}{\cftdotsep}
\renewcommand*{\cftdotsep}{1}
\renewcommand*{\cftchapterleader}{\cftdotfill{\cftchapterdotsep}}
\maxtocdepth{subsection}

%%% Выравнивание и переносы %%%
\tolerance 1414
\hbadness 1414
\emergencystretch 1.5em
\hfuzz 0.3pt
\vfuzz \hfuzz
\clubpenalty=10000
\widowpenalty=10000
\brokenpenalty=4991

%%% Кириллица в перечислениях %%%
\makeatletter
    \def\russian@Alph#1{\ifcase#1\or
       А\or Б\or В\or Г\or Д\or Е\or Ж\or
       И\or К\or Л\or М\or Н\or
       П\or Р\or С\or Т\or У\or Ф\or Х\or
       Ц\or Ш\or Щ\or Э\or Ю\or Я\else\xpg@ill@value{#1}{russian@Alph}\fi}
    \def\russian@alph#1{\ifcase#1\or
       а\or б\or в\or г\or д\or е\or ж\or
       и\or к\or л\or м\or н\or
       п\or р\or с\or т\or у\or ф\or х\or
       ц\or ш\or щ\or э\or ю\or я\else\xpg@ill@value{#1}{russian@alph}\fi}
\makeatother

%%% Рисунки и таблицы %%%
\usepackage{graphicx, caption, subcaption}
\graphicspath{{images/}}
\captionsetup[figure]{font=small, width=\textwidth, name=Рисунок, justification=centering}
\captionsetup[subfigure]{font=small}
\captionsetup[table]{singlelinecheck=false,font=small,width=\textwidth,justification=justified}
\captiondelim{ --- }
\setkeys{Gin}{width=\textwidth}
\renewcommand{\thesubfigure}{\asbuk{subfigure}}
\usepackage[section]{placeins}

%%% Гиперссылки %%%
\usepackage{hyperref}
\hypersetup{
    colorlinks=true,
    linktoc=all,
    linktocpage=true,
    linkcolor=red,
    citecolor=red
}

%%% Списки %%%
\usepackage{enumitem}
\renewcommand*{\labelitemi}{\normalfont{--}}
\makeatletter
    \AddEnumerateCounter{\asbuk}{\russian@alph}
\makeatother
\renewcommand{\labelenumii}{\asbuk{enumii})}
\renewcommand{\labelenumiii}{\arabic{enumiii})}
\setlist{noitemsep, leftmargin=*}
\setlist[1]{labelindent=\parindent}
\setlist[2]{leftmargin=\parindent}
\setlist[3]{leftmargin=\parindent}

%%% Счётчики %%%
\counterwithout{figure}{chapter}
\counterwithout{equation}{chapter}
\counterwithout{table}{chapter}

%%% Блоки кода (листинги) %%%
\usepackage{listings}
\usepackage{xcolor}
\lstset{
    basicstyle=\small\ttfamily,
    breaklines=true,
    frame=single,
    tabsize=2,
    showstringspaces=false,
    extendedchars=true,
    inputencoding=utf8
}

%%% Библиография %%%
\usepackage{csquotes}
\usepackage[%
backend=biber,
bibencoding=utf8,
sorting=none,
style=gost-numeric,
language=auto,
autolang=other,
sortcites=true,
movenames=false,
maxnames=5,
minnames=3,
doi=false,
isbn=false,
]{biblatex}[2016/09/17]
\DeclareDelimFormat{bibinitdelim}{}
\addbibresource{bibl.bib}

\DeclareSourcemap{
  \maps[datatype=bibtex]{
    \map{
        \step[fieldsource=title, match=\regexp{^\P{Cyrillic}*\p{Cyrillic}.*}, final]
        \step[fieldset=langid, fieldvalue={russian}]
    }
    \map{
        \step[fieldset=langid, fieldvalue={english}]
    }
  }
}

%%% Доп. пакеты %%%
\usepackage{longtable,ltcaption}
\usepackage{multirow,makecell}
\usepackage{booktabs}
\usepackage{soulutf8}
\usepackage{icomma}
\usepackage{hyphenat}
\usepackage{textcomp}
\usepackage{amsmath}

%%% Документ %%%
\begin{document}

\thispagestyle{empty}
\begin{center}
\textbf{МИНИСТЕРСТВО ЦИФРОВОГО РАЗВИТИЯ, СВЯЗИ}\\
\textbf{И МАССОВЫХ КОММУНИКАЦИЙ РОССИЙСКОЙ ФЕДЕРАЦИИ}

\vspace{14pt}

Федеральное государственное бюджетное образовательное учреждение\\
высшего образования\\
«Сибирский государственный университет телекоммуникаций и информатики»

\vspace{14pt}

Кафедра телекоммуникационных систем и вычислительных средств (ТС и ВС)
\end{center}

\vfill

\begin{center}
\textbf{РЕФЕРАТ}\\
по дисциплине\\
\textit{\guillemotleft Web-\foreignlanguage{russian}{технологии}\guillemotright}

\vspace{14pt}

по теме:\\
{\large\textbf{DNS-сервер BIND9, зоны forward/reverse,\\
динамическое обновление DNS (DDNS)}}
\end{center}

\vfill

\noindentСтудент:\\
\textit{Группа \textnumero~ИКС-531 \hfill Д.А.~Язиков}

\vspace{14pt}

\noindentПреподаватель:\\
\textit{ст.~преподаватель \hfill А.В.~Андреев}

\vfill

\begin{center}
Новосибирск 2025~г.
\end{center}
\newpage


\newpage
\setcounter{page}{2}
\OnehalfSpacing*

\tableofcontents*

\section*{Введение}
\addcontentsline{toc}{section}{Введение}

Система доменных имён (DNS) является одним из фундаментальных компонентов современных компьютерных сетей. Она обеспечивает преобразование символьных доменных имён в числовые IP-адреса. Без DNS пользователь был бы вынужден запоминать IP-адреса каждого сетевого ресурса вместо привычных имён.

В рамках данной практической работы рассматривается развёртывание собственного DNS-сервера на базе пакета \textbf{BIND9} (Berkeley Internet Name Domain) в операционной системе Ubuntu~20.04~LTS~\cite{bind9_admin}.

Цель работы~--- настройка полноценного DNS-сервера для локальной сети, включающая:
\begin{itemize}
  \item создание прямой зоны (forward zone) для преобразования имён в IP-адреса;
  \item создание обратной зоны (reverse zone) для преобразования IP-адресов в имена;
  \item настройку DDNS для авторегистрации клиентов через DHCP.
\end{itemize}

Стенд: виртуальная машина \texttt{gateway} (VirtualBox), интерфейсы \texttt{enp0s3}~(WAN/NAT) и \texttt{enp0s8}~(LAN, \texttt{192.168.N.1/24}).

\chapter{Теоретические сведения}
\label{ch:theory}

\section{Система доменных имён DNS}

DNS (Domain Name System) — иерархическая распределённая система именования узлов сети \cite{rfc1034}. Пространство имён DNS имеет древовидную структуру: в корне находится безымянный корневой домен, ниже расположены домены верхнего уровня (TLD: \texttt{.ru}, \texttt{.com}, \texttt{.local} и др.), затем домены второго уровня и так далее.

Основные понятия DNS:
\begin{itemize}
    \item \textbf{Зона} --- административная единица пространства имён, за которую отвечает конкретный DNS-сервер;
    \item \textbf{Авторитативный сервер} (authoritative) --- сервер, хранящий оригинальные записи зоны и отвечающий на запросы с пометкой «авторитативный ответ»;
    \item \textbf{Рекурсивный резолвер} --- сервер, выполняющий цепочку запросов к другим DNS-серверам от имени клиента;
    \item \textbf{Форвардер} (forwarder) --- сервер, которому локальный DNS-сервер перенаправляет запросы для не локальных зон.
\end{itemize}

\section{Типы ресурсных записей DNS}

Каждая запись в зоне DNS имеет определённый тип. Основные типы, используемые в данной работе, приведены в таблице~\ref{tab:rr}.

\begin{table}[h]
    \caption{Основные типы ресурсных записей DNS}
    \label{tab:rr}
    \begin{tabular}{|l|p{10cm}|}
        \hline
        \textbf{Тип} & \textbf{Назначение} \\
        \hline
        SOA & Start of Authority --- описывает зону: первичный NS, e-mail администратора, параметры обновления \\
        \hline
        NS  & Name Server --- указывает авторитативный сервер зоны \\
        \hline
        A   & Адресная запись --- сопоставляет имя хоста с IPv4-адресом \\
        \hline
        PTR & Pointer --- обратная запись, сопоставляет IP-адрес с именем хоста \\
        \hline
        CNAME & Canonical Name --- псевдоним для другого доменного имени \\
        \hline
        MX  & Mail Exchanger --- указывает почтовый сервер домена \\
        \hline
    \end{tabular}
\end{table}

\section{Прямая и обратная зоны}

\textbf{Прямая зона} (forward zone) используется для разрешения имён в IP-адреса. Имя зоны совпадает с именем домена, например \texttt{STUDENT.GROUP.local}. Файл зоны содержит записи типов SOA, NS, A.

\textbf{Обратная зона} (reverse zone) используется для разрешения IP-адресов в имена. Имя зоны строится по принципу инвертирования первых октетов подсети с добавлением суффикса \texttt{.in-addr.arpa}. Например, для сети \texttt{192.168.N.0/24} зона будет называться \texttt{N.168.192.in-addr.arpa}. Файл зоны содержит записи типов SOA, NS, PTR.

\section{BIND9}

BIND (Berkeley Internet Name Domain) версии 9 --- наиболее распространённое и функционально полное ПО DNS-сервера \cite{bind9_admin}. В Ubuntu 20.04 устанавливается из стандартного репозитория командой \texttt{apt install bind9}. Конфигурационные файлы располагаются в каталоге \texttt{/etc/bind/}:
\begin{itemize}
    \item \texttt{named.conf} --- главный файл, включает остальные;
    \item \texttt{named.conf.options} --- глобальные настройки (форвардеры, listen-on, dnssec);
    \item \texttt{named.conf.local} --- описание локальных зон;
    \item файлы зон хранятся в \texttt{/var/lib/bind/} --- в этом каталоге BIND имеет права на запись, что необходимо для DDNS.
\end{itemize}

\section{Динамическое обновление DNS (DDNS)}

DDNS (Dynamic DNS, RFC~2136) позволяет клиентам или другим серверам автоматически обновлять записи DNS в зоне без ручного вмешательства администратора \cite{rfc2136}. В данной работе DDNS используется совместно с DHCP-сервером \texttt{isc-dhcp-server}: при получении клиентом IP-адреса DHCP-сервер автоматически создаёт A-запись и PTR-запись в DNS.

Для аутентификации обновлений используется TSIG-ключ (Transaction SIGnature) --- симметричный ключ HMAC-MD5, генерируемый утилитой \texttt{rndc-confgen}. Ключ хранится в файле \texttt{/etc/bind/rndc.key} и должен быть доступен как BIND9, так и DHCP-серверу.

\endinput

\chapter{Описание стенда и подготовка}
\label{ch:stand}

\section{Конфигурация виртуального стенда}

Лабораторная работа выполняется на виртуальной машине \texttt{gateway} под управлением ОС \textbf{Ubuntu 20.04.5 LTS (Focal Fossa)}, развёрнутой в среде Oracle VirtualBox. ВМ имеет два сетевых адаптера:
\begin{itemize}
    \item \texttt{enp0s3} --- WAN-интерфейс, подключён в режиме NAT; получает адрес \texttt{10.0.2.15/24}, шлюз \texttt{10.0.2.2};
    \item \texttt{enp0s8} --- LAN-интерфейс, подключён к внутренней сети (Internal Network); статический адрес \texttt{192.168.N.1/24}.
\end{itemize}

Лаба 5 выполняется на той же ВМ, что и лаба 4. К моменту начала лабы 5 уже настроены: маршрутизация (ip\_forward + iptables MASQUERADE) и DHCP-сервер \texttt{isc-dhcp-server}.

\section{Предварительная подготовка}

Перед установкой BIND9 необходимо выполнить ряд подготовительных шагов.

\subsection{Установка hostname}

Имя хоста сервера задаётся командой:
\begin{lstlisting}
hostnamectl set-hostname SERVERHOSTNAME
\end{lstlisting}
где \texttt{SERVERHOSTNAME} --- выбранное имя ВМ (например, \texttt{gateway}). Файл \texttt{/etc/hosts} редактируется вручную:
\begin{lstlisting}
nano /etc/hosts
# добавить строку:
127.0.1.1   SERVERHOSTNAME
\end{lstlisting}

\subsection{Удаление DNS DNAT-правил}

В лабе 4 в файл \texttt{/etc/iptables/rules.v4} были добавлены правила перенаправления DNS-запросов на Google~8.8.8.8. Теперь эти правила необходимо удалить, поскольку DNS-сервер будет обрабатывать запросы самостоятельно:
\begin{lstlisting}
nano /etc/iptables/rules.v4
# удалить строки:
-A PREROUTING -i enp0s8 -p tcp -m tcp --dport 53 \
  -j DNAT --to-destination 8.8.8.8:53
-A PREROUTING -i enp0s8 -p udp -m udp --dport 53 \
  -j DNAT --to-destination 8.8.8.8:53
\end{lstlisting}

После редактирования применить изменения и перезагрузить ВМ:
\begin{lstlisting}
reboot
\end{lstlisting}

\section{Установка BIND9}

Установка пакетов производится стандартными средствами менеджера пакетов:
\begin{lstlisting}
apt-get update
apt install bind9
apt install dnsutils
\end{lstlisting}

Пакет \texttt{dnsutils} предоставляет утилиты \texttt{nslookup} и \texttt{dig}, необходимые для проверки работы DNS.

\endinput

\section{Настройка DNS-сервера BIND9}

\subsection{Настройка named.conf.options}

\begin{lstlisting}
nano /etc/bind/named.conf.options
\end{lstlisting}

\begin{lstlisting}
options {
    directory "/var/cache/bind";
    forwarders { 8.8.8.8; };
    dnssec-validation auto;
    auth-nxdomain no;
    listen-on { 127.0.0.1; 192.168.N.1; };
};
\end{lstlisting}

\texttt{forwarders}~--- DNS Google для внешних зон; \texttt{listen-on}~--- принимать запросы только с localhost и LAN-интерфейса.

\subsection{Настройка named.conf.local}

\begin{lstlisting}
nano /etc/bind/named.conf.local
\end{lstlisting}

\begin{lstlisting}
include "/etc/bind/rndc.key";
controls {
    inet 127.0.0.1 allow { localhost; } keys { rndc-key; };
};
zone "STUDENT.GROUP.local" IN {
    type master;
    file "/var/lib/bind/forward.db";
    allow-update { key rndc-key; };
};
zone "N.168.192.in-addr.arpa" IN {
    type master;
    file "/var/lib/bind/reverse.db";
    allow-update { key rndc-key; };
};
\end{lstlisting}

Файлы зон в \texttt{/var/lib/bind/}: здесь BIND имеет права на запись, что нужно для DDNS (журнальные файлы \texttt{.jnl}).

\subsection{Файл прямой зоны forward.db}

\begin{lstlisting}
nano /var/lib/bind/forward.db
\end{lstlisting}

\begin{lstlisting}
$TTL 86400
STUDENT.GROUP.local. IN SOA SERVERHOSTNAME.STUDENT.GROUP.local.
    admin.STUDENT.GROUP.local. (
        20110103 ; Serial
        604800   ; Refresh
        86400    ; Retry
        2419200  ; Expire
        604800 ) ; Negative Cache TTL
@       IN NS  SERVERHOSTNAME.STUDENT.GROUP.local.
@       IN A   192.168.N.1
localhost IN A  127.0.0.1
SERVERHOSTNAME IN A 192.168.N.1
\end{lstlisting}

\subsection{Файл обратной зоны reverse.db}

\begin{lstlisting}
nano /var/lib/bind/reverse.db
\end{lstlisting}

\begin{lstlisting}
$TTL 86400
N.168.192.in-addr.arpa. IN SOA STUDENT.GROUP.local.
    STUDENT.GROUP.local. (
        20110104 ; Serial
        10800    ; Refresh
        3600     ; Retry
        604800   ; Expire
        3600 )   ; Minimum
        IN NS  SERVERHOSTNAME.STUDENT.GROUP.local.
1       IN PTR STUDENT.GROUP.local.
1       IN PTR SERVERHOSTNAME.STUDENT.GROUP.local.
\end{lstlisting}

\subsection{Проверка и перезапуск BIND9}

\begin{lstlisting}
named-checkconf
named-checkzone STUDENT.GROUP.local /var/lib/bind/forward.db
named-checkzone N.168.192.in-addr.arpa /var/lib/bind/reverse.db
systemctl restart bind9
\end{lstlisting}

\subsection{Настройка netplan}

\begin{lstlisting}
nano /etc/netplan/00-installer-config.yaml
\end{lstlisting}

\begin{lstlisting}
network:
  ethernets:
    enp0s3:
      dhcp4: no
      addresses: [10.0.2.15/24]
      gateway4: 10.0.2.2
    enp0s8:
      dhcp4: no
      addresses: [192.168.N.1/24]
      nameservers:
        addresses: [192.168.N.1]
        search: [STUDENT.GROUP.local]
  version: 2
\end{lstlisting}

\begin{lstlisting}
netplan apply && reboot
\end{lstlisting}

\begin{lstlisting}
nslookup SERVERHOSTNAME
nslookup 192.168.N.1
\end{lstlisting}

\subsection{Настройка DDNS}

\textbf{Копирование ключа RNDC:}
\begin{lstlisting}
cp -Rr /etc/bind/rndc.key /etc/dhcp/ddns-keys/
\end{lstlisting}

\textbf{Обновление dhcpd.conf:}
\begin{lstlisting}
include "/etc/dhcp/ddns-keys/rndc.key";
ddns-updates on;
ddns-update-style standard;
ddns-domainname "STUDENT.GROUP.local";
zone STUDENT.GROUP.local. {
    primary 192.168.N.1; key rndc-key;
}
zone N.168.192.in-addr.arpa. {
    primary 192.168.N.1; key rndc-key;
}
authoritative;
subnet 192.168.N.0 netmask 255.255.255.0 {
    range 192.168.N.10 192.168.N.254;
    option domain-name-servers 192.168.N.1;
    option domain-name "STUDENT.GROUP.local";
    option routers 192.168.N.1;
    option broadcast-address 192.168.N.255;
    default-lease-time 604800;
    max-lease-time 604800;
}
\end{lstlisting}

\textbf{Перезапуск и проверка:}
\begin{lstlisting}
systemctl restart bind9
service isc-dhcp-server restart
nslookup DESKTOPNAME01
tail -40 /var/log/syslog
\end{lstlisting}

\section*{Заключение}
\addcontentsline{toc}{section}{Заключение}

В ходе выполнения практической работы был развёрнут и настроен DNS-сервер BIND9 на виртуальной машине \texttt{gateway} под управлением Ubuntu~20.04~LTS. Созданы прямая зона \texttt{STUDENT.GROUP.local} и обратная зона \texttt{N.168.192.in-addr.arpa}.

Настройка DDNS с ключом \texttt{rndc-key} позволила автоматизировать регистрацию клиентов в DNS при получении IP-адреса через DHCP. Клиентские машины автоматически получают доменные имена \texttt{HOSTNAME.STUDENT.GROUP.local}.

Приобретённые навыки~--- администрирование BIND9, работа с файлами зон, настройка DDNS~--- являются базовыми компетенциями сетевого администратора.


\printbibliography[title=Список использованных источников]

\end{document}
